% ------------------------------------------------------------
% Page 103
% ------------------------------------------------------------

Il prit congé avec les membres de la Commission, mais, en se retournant sur le pas de la porte
il ajouta :
« Meilleurs vœux pour votre mariage ; trois jours c’est fort peu, alors prenez-en quatre ou
cinq si c’est nécessaire.
- Merci Monsieur l’Inspecteur. je n’abuserai pas ».

Ainsi se termina cette visite dont je garde un souvenir très précis ; j’ai réalisé plus tard que cet
homme avait certainement perçu en moi un certain désarroi, et qu’il s’était attaché, tout au
long de l’entretien, à me réconforter. Je n’ai plus revu cet inspecteur muté quelques mois plus
tard.

\textbf{Mariage et arrivée de ma jeune femme.}

Le remplacement de la fameuse fenêtre m’a rendu plus serein. Et voilà que dans les derniers
jours d’octobre, les visites quasi simultanées de ma famille et de ma belle-famille mettent fin
aux jours d’inquiétude et de désarroi.

Ma mère arrive avec JA, voiturier du Pompidou, avec tables, chaises, matelas, provisions de
première nécessité, et même une suspension qu’elle me « prête » en attendant. Et c’est ma
belle famille qui apporte une belle cuisinière en fonte émaillée, un petit buffet en bois blanc
(cadeau de la grand-mère), une penderie ouvragée par mon beau-père, le tout accompagnant
cadeaux de mariage et trousseau méticuleusement conçu et préparé pour l’unique fille de la
maison.

Pas d’armoire, mais où l’aurions-nous logée ? Je pense alors à la grande armoire du fond de la
classe, dont le bas « est plein de vide » : ce sera suffisant pour loger lingerie, draps et
couvertures.
Je n’oublie pas le vélo remis en état, ni la canne à pêche, jugée indispensable par mon beau-
père pour quelqu’un qui vit aussi près d’une belle rivière.
Comme par enchantement, la petite sacristie, bien suffisamment meublée, devient un petit
« deux pièces » relativement accueillant, et c’est avec un moral tout neuf que j’envisage
l’avenir.
Fini ce mois d’octobre, mais pas oublié !! Il revient souvent dans nos conversations. J’ai été
très, très long, mais pouvais-je être bref ? J’ai à peine parlé de mes élèves, et c’est un comble
pour un « instit » !

\textbf{Mes élèves}

Après un mois fort agité, passé à l’école-temple des Oubrets, il est normal, fonction oblige,
que l’instit en titre consacre quelques lignes à ses élèves. Ils n’étaient que trois :

\underline{Martin} est plein de bonne volonté, pour me faire plaisir, mais je ne le sens pas très motivé. Je
ne l’ai pas présenté au CEP ; Peut-être ai-je eu tort ? après avoir participé à plusieurs
commissions du C.E.P., j’ai réalisé qu’il avait des chances. Il aime les moutons et me tient au
courant de l’évolution de son troupeau, de ses besogne... Il réussira fort bien, quelques années
plus tard, comme responsable d’un grand troupeau laitier sur le Causse Méjean.
